\documentclass[a4paper,12pt,oneside]{report}

% ------------------------------------------------
% LANGUAGE & FONTS
% ------------------------------------------------
\usepackage[british]{babel}
\usepackage[T1]{fontenc}
\usepackage[utf8]{inputenc}
\usepackage{newtxtext,newtxmath} % Times New Roman-like fonts for pdflatex

% ------------------------------------------------
% PAGE GEOMETRY
% ------------------------------------------------
\usepackage[a4paper,
    left=3.5cm,
    right=1.25cm,
    top=2.5cm,
    bottom=1.25cm
]{geometry}

% ------------------------------------------------
% LINE SPACING & INDENTATION
% ------------------------------------------------
\usepackage{setspace}
\doublespacing
\setlength{\parindent}{1.25cm}

% ------------------------------------------------
% CHAPTER, SECTION, SUBSECTION FORMATTING
% ------------------------------------------------
\usepackage{titlesec}
% Chapter Title - 16pt, Bold, Capital, Centered
\titleformat{\chapter}[block]
  {\centering\normalfont\bfseries\fontsize{16}{20}\selectfont}
  {\thechapter.}{1em}{\MakeUppercase}

% Section - 14pt, Bold
\titleformat{\section}
  {\normalfont\bfseries\fontsize{14}{16}\selectfont}{\thesection}{1em}{}

% Subsection - 12pt, Bold
\titleformat{\subsection}
  {\normalfont\bfseries\fontsize{12}{14}\selectfont}{\thesubsection}{1em}{}

% ------------------------------------------------
% CAPTION SETTINGS
% ------------------------------------------------
\usepackage{caption}
\captionsetup[figure]{
    font={it,small},
    labelfont={bf},
    labelsep=period,
    justification=centering,
    name=Figure
}
\captionsetup[table]{
    font={it,small},
    labelfont={bf},
    labelsep=period,
    justification=centering,
    name=Table,
    position=top
}

% ------------------------------------------------
% NUMBERING FORMAT
% ------------------------------------------------
\numberwithin{equation}{chapter}
\numberwithin{figure}{chapter}
\numberwithin{table}{chapter}

% ------------------------------------------------
% OTHER PACKAGES
% ------------------------------------------------
\usepackage{graphicx}
\usepackage{amsmath}
\usepackage{listings}
\usepackage{float}
\usepackage{makeidx}
\usepackage{hyperref}
\usepackage{xcolor}
\usepackage{microtype}

\lstset{
    basicstyle=\ttfamily\footnotesize,
    breaklines=true
}

% ------------------------------------------------
% TITLE PAGE
% ------------------------------------------------
\begin{document}

\begin{titlepage}
\begin{center}
    \Large\textbf{FOUR WEEK TRAINING REPORT}\\[3mm]
    \large{at}\\[3mm]
    \Large\textbf{Academic Advancement of Information Technology, Mohali}\\[6mm]
    \normalsize{SUBMITTED IN PARTIAL FULFILLMENT OF THE REQUIREMENTS FOR THE AWARD OF DEGREE OF}\\[3mm]
    \Large\textbf{BACHELOR OF TECHNOLOGY}\\[3mm]
    \large{in Computer Science and Engineering}\\[9mm]
    \includegraphics[height=6cm]{gndeclogo.png}\\[6mm]
    \large{JUNE–JULY 2025}\\[9mm]
    \textbf{SUBMITTED BY:}\\
    NAME: Ayush Mehta\\
    UNIVERSITY ROLL NO.: 2302489
\end{center}

\vspace{10mm}
\begin{center}
    \textbf{Department of Computer Science and Engineering}\\
    \textbf{Guru Nanak Dev Engineering College}\\
    \textbf{Ludhiana, 141006}
\end{center}
\end{titlepage}

% ------------------------------------------------
% CERTIFICATE PAGE (TO BE INSERTED FROM COMPANY)
% ------------------------------------------------
\newpage
\begin{center}
    \large\textbf{CERTIFICATE}\\[6mm]
    \includegraphics[width=\textwidth, keepaspectratio]{certificate.jpg}
\end{center}

% ------------------------------------------------
% CANDIDATE’S DECLARATION
% ------------------------------------------------
\newpage
\begin{center}
    \large\textbf{CANDIDATE'S DECLARATION}
\end{center}

I, \textbf{Ayush Mehta}, hereby declare that I have undertaken four-week Web Development training from \textbf{Academic Advancement of Information Technology, Mohali} during the period from 26 June 2025 to 26 July 2025 in partial fulfillment of the requirements for the award of the degree of \textbf{B.Tech. (Computer Science and Engineering)} at \textbf{Guru Nanak Dev Engineering College, Ludhiana}. The work presented in this training report is an authentic record of my training.

\vspace{15mm}
\noindent
\textbf{(Ayush Mehta)}\\
Roll No.: 2302489\\[10mm]
The four week industrial training Viva--Voce Examination of \rule{4cm}{0.4pt} has been held on \rule{3cm}{0.4pt} and accepted.

\vspace{45mm}

\noindent
\begin{tabular}{p{0.45\textwidth}p{0.45\textwidth}}
\centering\textbf{Signature of External Examiner} & \centering\textbf{Signature of Internal Examiner} \\
\end{tabular}

% ------------------------------------------------
% ABSTRACT
% ------------------------------------------------
\newpage
\begin{center}
    \large\textbf{ABSTRACT}
\end{center}

This report summarizes the four-week industrial training in Web Development undertaken at \textbf{Academic Advancement of Information Technology (A2IT), Mohali}. The training primarily focused on learning the fundamentals of front-end web technologies, including \textbf{HTML} and \textbf{CSS}, along with an introductory understanding of \textbf{JavaScript}. 

As a beginner to web development, this training provided me with a strong foundation in creating structured, styled, and responsive web pages. The sessions covered essential concepts of website design and layout, enabling me to understand how the various components of a web application interact. 

Towards the end of the training, I developed a small project—a \textbf{Scientific Web Calculator}—which allowed me to apply the knowledge gained during the sessions. Although simple, this project served as a practical exercise to consolidate the learning outcomes. Overall, the training proved to be an invaluable starting point for my journey into web development and helped me build confidence in working with core web technologies.

% ------------------------------------------------
% ACKNOWLEDGEMENT, CONTENTS, ETC.
% ------------------------------------------------

\newpage
\begin{center}
    \large\textbf{ACKNOWLEDGEMENT}
\end{center}

I express my deepest sense of gratitude to \textbf{Dr.~Sehijpal Singh}, Principal, Guru Nanak Dev Engineering College, Ludhiana, for providing the necessary facilities and environment for carrying out the training successfully. I am equally thankful to \textbf{Dr.~Kiran Jyoti}, Head, Department of Computer Science and Engineering, for her valuable support, motivation, and guidance during the course of the training.

I am sincerely thankful to \textbf{Mr.~Jaswant Singh} and \textbf{Ms.~Kuljit Kaur}, Training Coordinators, Department of Computer Science and Engineering, for their constant guidance, encouragement, and for providing valuable instructions regarding the preparation of this report and the training documentation.

I would also like to express my heartfelt thanks to the management and staff of \textbf{Academic Advancement of Information Technology (A2IT), Mohali} for providing me the opportunity to undergo industrial training. I extend my sincere appreciation to \textbf{Ms.~Payal Karn} for her continuous guidance, insightful lectures, and support throughout the training period. I am also thankful to \textbf{Mr.~Rajeev}, Director of Technology, A2IT, for facilitating the overall training program, and to \textbf{Ms.~Harpreet Kaur}, Senior Manager (Human Resources), A2IT, for her assistance in administrative and certification matters.

Lastly, I extend my gratitude to all my faculty members, friends, and family who directly or indirectly helped me during this training and in preparing this report.

\vspace{10mm}
\noindent
\textbf{(Ayush Mehta)}\\
Roll No.: 2302489

\tableofcontents
\listoffigures
\listoftables

% ------------------------------------------------
% CHAPTERS
% ------------------------------------------------
\newpage
\chapter{INTRODUCTION}

\section{BACKGROUND}

The field of web development has rapidly evolved over the past few decades, becoming one of the most dynamic and essential domains in computer science and information technology. The internet, once a medium for static information display, has transformed into a highly interactive platform enabling businesses, educational institutions, and individuals to connect globally.

In the early stages of web design, developers primarily relied on simple \textbf{HTML (HyperText Markup Language)} to create basic webpages. As the need for visual appeal and structure grew, \textbf{CSS (Cascading Style Sheets)} was introduced, allowing developers to separate content from design and apply consistent styling across webpages. Over time, the demand for interactivity gave rise to \textbf{JavaScript}, which added dynamic features such as form validation, animations, and responsive user interfaces.

Today, web development encompasses a wide variety of technologies and frameworks designed to enhance both performance and user experience. However, understanding the fundamental building blocks—HTML, CSS, and JavaScript—remains crucial for every aspiring web developer. Mastery of these core technologies lays the foundation for advanced front-end and full-stack development. 

This training program was designed with beginners in mind, focusing on practical, hands-on exposure to these core technologies. As a newcomer to web development, this four-week program provided me the opportunity to move from zero prior experience to confidently developing simple, well-structured, and visually appealing webpages.

\section{OBJECTIVE OF THE TRAINING}

The primary objective of the four-week industrial training in Web Development was to equip participants with essential skills and foundational understanding of front-end technologies. The training aimed to foster practical learning through continuous implementation rather than purely theoretical study.

The key objectives of the training were as follows:

\begin{itemize}
    \item To understand and implement semantic HTML for structured and accessible web pages.
    \item To learn CSS fundamentals for styling, layout design, and responsive web development.
    \item To gain familiarity with JavaScript basics including variables, data types, operators, functions, conditionals, loops, and events.
    \item To understand the Document Object Model (DOM) and how JavaScript interacts with HTML elements.
    \item To design responsive and accessible web interfaces using modern CSS techniques and frameworks.
    \item To develop small-scale, functional web projects demonstrating integration of HTML, CSS, and JavaScript.
    \item To build a simple \textbf{Scientific Web Calculator} as a final project, consolidating the concepts learned throughout the training.
\end{itemize}

The training was conducted in a structured, progressive manner—beginning with the basics of markup and styling and culminating in interactive page design and scripting. This approach helped bridge the gap between conceptual understanding and practical application.

\section{OVERVIEW OF WEB DEVELOPMENT TRAINING}

The web development training covered the three core technologies that form the backbone of front-end development:

\begin{itemize}
    \item \textbf{HTML (HyperText Markup Language):} Used to structure web content with elements like headings, paragraphs, tables, forms, images, and multimedia.
    \item \textbf{CSS (Cascading Style Sheets):} Used to style and visually enhance HTML elements, enabling layout control, color schemes, typography, and responsive design.
    \item \textbf{JavaScript:} A lightweight scripting language used to add interactivity, handle user inputs, and manipulate the Document Object Model dynamically.
\end{itemize}

Throughout the training, emphasis was placed on clean, semantic coding practices and the use of external style sheets for maintainability. Responsive design techniques, accessibility considerations, and code validation were also highlighted. 

In the final phase of training, the concepts were integrated through the creation of a \textbf{Scientific Web Calculator} project. This project demonstrated the application of HTML structure, CSS styling, and JavaScript functionality to create a practical and interactive web-based tool.

\section{IMPORTANCE OF WEB DEVELOPMENT IN THE MODERN ERA}

Web development is one of the most sought-after skills in the modern digital age. Almost every organization—from startups to global enterprises—relies on web applications to deliver products, services, and information to users worldwide. Understanding the fundamentals of web technologies is essential not only for computer science students but also for anyone interested in digital innovation.

Some of the key reasons for the importance of web development include:

\begin{itemize}
    \item \textbf{Universal Accessibility:} Websites serve as globally accessible platforms for information, communication, and commerce.
    \item \textbf{Career Relevance:} Proficiency in HTML, CSS, and JavaScript is a foundational skill set required for many modern software roles.
    \item \textbf{Creative and Analytical Balance:} Web development uniquely combines logical problem-solving with visual and creative design.
    \item \textbf{Scalability and Flexibility:} Websites and web apps can be scaled easily across devices and platforms, making them cost-effective and efficient.
    \item \textbf{Continuous Growth:} The web ecosystem is constantly evolving, with new tools, frameworks, and best practices emerging regularly.
\end{itemize}

Through this training, I have gained a beginner-level yet substantial understanding of how web pages are structured, styled, and made interactive. This foundation paves the way for future exploration into advanced frameworks and backend technologies.

\section{SCOPE OF TRAINING}

The training covered both theoretical and practical aspects of front-end development, emphasizing implementation-based learning. The scope of work and skill development can be summarized as follows:

\textbf{HTML and CSS Development:}
\begin{itemize}
    \item Creation of semantic HTML structures including tables, lists, and forms.
    \item Integration of multimedia elements such as images, audio, and video.
    \item Application of CSS for page layout, typography, spacing, and visual aesthetics.
    \item Implementation of responsive design principles using relative units and media queries.
    \item Understanding the box model, selectors, and cascading hierarchy.
\end{itemize}

\textbf{JavaScript Fundamentals:}
\begin{itemize}
    \item Learning basic syntax, data types, and operators.
    \item Using control flow structures like conditionals and loops.
    \item Understanding functions, events, and DOM manipulation.
    \item Developing small scripts for user interaction and dynamic page behavior.
\end{itemize}

\textbf{Final Project:}
\begin{itemize}
    \item Designing and developing a simple \textbf{Scientific Web Calculator}.
    \item Implementing user interaction through event handling.
    \item Managing form inputs, mathematical operations, and display updates via JavaScript.
    \item Styling the interface using CSS to ensure readability and usability.
\end{itemize}

Overall, the training provided a strong introduction to web development principles and practices, enabling me to create structured, styled, and functional web applications independently.


\newpage
\begin{center}
    \Large\textbf{CHAPTER 2}\\[2mm]
    \large\textbf{TRAINING WORK UNDERTAKEN}
\end{center}

The four-week industrial training in Web Development at \textbf{Academic Advancement of Information Technology (A2IT), Mohali} was aimed at introducing the fundamental concepts and practices of modern web design and development. The training was primarily focused on the core building blocks of front-end technologies—\textbf{HTML}, \textbf{CSS}, and the introductory concepts of \textbf{JavaScript}. Over the course of the training, emphasis was placed on understanding the logical structure of web pages, styling principles, and the integration of interactivity to create a complete and responsive web experience.

\vspace{5mm}
\noindent
\textbf{2.1 WEEK 1 – INTRODUCTION TO HTML AND WEB STRUCTURE}

The first week of the training concentrated on learning the fundamental concepts of web page creation using \textbf{HTML (HyperText Markup Language)}. The objective was to understand the structure of web documents and how browsers interpret markup elements to display content effectively.

\textbf{Topics Covered:}
\begin{itemize}
    \item Introduction to HTML and its role in web development.
    \item Document structure using \texttt{<!DOCTYPE html>}, \texttt{<html>}, \texttt{<head>}, and \texttt{<body>} tags.
    \item Use of text formatting, paragraph, heading, and list elements.
    \item Embedding images, hyperlinks, and multimedia content.
    \item Overview of semantic HTML elements such as \texttt{<header>}, \texttt{<section>}, \texttt{<article>}, and \texttt{<footer>}.
\end{itemize}

By the end of this stage, participants developed the ability to construct multi-page static websites with proper structure and hierarchy, forming the basis for all subsequent design and scripting activities.

\newpage
\noindent
\textbf{2.2 WEEK 2 – INTRODUCTION TO CSS AND PAGE STYLING}

The second week was dedicated to \textbf{CSS (Cascading Style Sheets)}, focusing on the visual and aesthetic aspects of web pages. Learners were trained in applying styles to HTML elements and developing visually coherent layouts that adhered to modern design standards.

\textbf{Topics Covered:}
\begin{itemize}
    \item Application of CSS through inline, internal, and external stylesheets.
    \item Text formatting, font properties, and color management.
    \item Understanding and applying the box model: margins, borders, padding, and content.
    \item Implementation of layouts using Flexbox and CSS Grid.
    \item Introduction to responsive web design and use of media queries.
\end{itemize}

Through practical sessions, participants achieved proficiency in organizing and styling webpage components systematically, ensuring both consistency and readability across devices.

\vspace{5mm}
\noindent
\textbf{2.3 WEEK 3 – INTRODUCTION TO JAVASCRIPT AND INTERACTIVITY}

The third week introduced the scripting component of front-end development using \textbf{JavaScript}. This segment emphasized the role of client-side scripting in enhancing interactivity and functionality within web applications.

\textbf{Topics Covered:}
\begin{itemize}
    \item Fundamentals of JavaScript syntax, variables, and data types.
    \item Implementation of conditional statements, loops, and functions.
    \item Interaction with HTML elements through the Document Object Model (DOM).
    \item Event handling and basic form validation.
    \item Dynamic content manipulation using JavaScript.
\end{itemize}

By the completion of this week, learners were capable of creating interactive web pages capable of responding to user inputs and executing simple client-side logic effectively.

\newpage
\noindent
\textbf{2.4 WEEK 4 – FINAL PROJECT DEVELOPMENT: SCIENTIFIC WEB CALCULATOR}

The final week was devoted to applying the knowledge gained during the earlier sessions in the form of a small-scale project. The chosen project was the development of a \textbf{Scientific Web Calculator}, designed to perform essential arithmetic and scientific operations through a user-friendly interface.

\textbf{Project Highlights:}
\begin{itemize}
    \item Interface design using HTML and CSS to ensure usability and readability.
    \item JavaScript-based logic for arithmetic and basic scientific calculations.
    \item Implementation of event handling for button operations and input validation.
    \item Layout optimization for responsiveness on varying screen sizes.
    \item Error handling for invalid or incomplete user inputs.
\end{itemize}

This project effectively consolidated the skills acquired throughout the training, providing hands-on experience in the end-to-end development of a functional web application.

\vspace{5mm}
\noindent
\textbf{2.5 TOOLS AND TECHNOLOGIES USED}

\begin{center}
\begin{tabular}{|p{5cm}|p{8cm}|}
\hline
\textbf{Technology / Tool} & \textbf{Purpose / Usage} \\
\hline
HTML5 & Structure and layout of web content \\
\hline
CSS3 & Styling and presentation of web pages \\
\hline
JavaScript & Adding logic and interactivity \\
\hline
Visual Studio Code & Code editing and project organization \\
\hline
Google Chrome Developer Tools & Debugging and layout inspection \\
\hline
Git and GitHub & Version control and repository management \\
\hline
\end{tabular}
\end{center}

\vspace{3mm}
The training concluded with a comprehensive understanding of web development fundamentals. Although introductory in nature, it provided a solid technical foundation for further exploration into advanced concepts such as responsive design frameworks, client-server communication, and backend development.


\newpage
\begin{center}
    \Large\textbf{CHAPTER 3}\\[2mm]
    \large\textbf{RESULTS AND DISCUSSIONS}
\end{center}

\noindent
\textbf{3.1 Overview of the Project Output}

The final project developed during the industrial training was a \textbf{Scientific Web Calculator} built using \textbf{HTML}, \textbf{CSS}, and \textbf{JavaScript}. The objective of the project was to design and implement a responsive and user-friendly calculator capable of performing both basic arithmetic and fundamental scientific operations such as trigonometric, logarithmic, exponential, and square root calculations.

The project serves as a practical implementation of front-end web development concepts learned during the training period. It focuses on dynamic user interaction, functional computation through JavaScript logic, and a visually appealing, responsive interface designed using modern CSS styling principles.

\bigskip
\noindent
\textbf{3.2 Implementation Results}

The Scientific Calculator was implemented as a modular, browser-based application. The development process was divided into three main stages—structural design, styling, and scripting—each focusing on a specific technological component.

\bigskip
\noindent
\textbf{3.2.1 Interface Layout and Design}

The calculator interface was created using \textbf{HTML5} to define the logical structure of the application. The layout consists of two primary sections:
\begin{itemize}
    \item \textbf{Main Calculator Grid:} Includes numeric buttons (0–9), arithmetic operations (+, –, ×, ÷), and special keys such as \textit{AC}, \textit{DEL}, and \textit{=}.
    \item \textbf{Scientific Function Panel:} Contains buttons for trigonometric functions (\textit{sin}, \textit{cos}, \textit{tan}), logarithmic and exponential functions (\textit{log}, \textit{exp}), square root, and mathematical constants such as \textit{$\pi$}.
\end{itemize}

This structure ensures that both standard and scientific computations can be performed from a single, organized interface.

\bigskip
\noindent
\newpage
\textbf{3.2.2 Visual Styling and Responsiveness}

The visual presentation of the calculator was developed using \textbf{CSS3}. The objective was to achieve a modern and intuitive appearance suitable for both desktop and mobile screens. 

\textbf{Key Features:}
\begin{itemize}
    \item \textbf{Dark Metallic Theme:} The background and button elements utilize dark gray and gradient tones to give a professional, high-contrast appearance.
    \item \textbf{Neon Green Display:} The calculator output screen is styled with neon green text on a dark background to simulate an LCD effect.
    \item \textbf{Hover and Click Effects:} Smooth transitions and shadow effects are used to indicate interactivity.
    \item \textbf{Responsive Layout:} The grid layout adapts to different screen sizes, ensuring usability across devices.
\end{itemize}

These design elements contribute to a polished user experience, maintaining both functionality and aesthetic appeal.

\bigskip
\noindent
\textbf{3.2.3 Functional Logic and Interactivity}

The computational logic was implemented using \textbf{JavaScript (ES6)}. The calculator operates through an object-oriented approach, encapsulated in a \texttt{Calculator} class responsible for handling user input, updating the display, and evaluating expressions.

\textbf{Core Features Implemented:}
\begin{itemize}
    \item \textbf{Input Handling:} Captures numeric, operational, and functional button presses using event listeners.
    \item \textbf{Expression Evaluation:} Mathematical expressions are preprocessed and evaluated using JavaScript’s \texttt{eval()} function with custom preprocessing for trigonometric and logarithmic operations.
    \item \textbf{Trigonometric Functions in Degrees:} Custom functions such as \texttt{sinDeg()}, \texttt{cosDeg()}, and \texttt{tanDeg()} ensure angle inputs are interpreted in degrees rather than radians.
    \item \textbf{Error Handling:} Displays “Error” for invalid expressions or undefined results to prevent unexpected behavior.
    \item \textbf{Dynamic Font Scaling:} The display font size adjusts automatically based on the length of the input or result.
\end{itemize}

This logic ensures smooth interaction, accurate computation, and a responsive user interface suitable for educational and demonstrative purposes.

\bigskip
\noindent
\textbf{3.2.4 Output Screens}

The following figures illustrate the functional output of the developed Scientific Calculator web application:

\begin{itemize}
    \item Figure 3.1 – Home screen layout of the Scientific Calculator interface.
    \item Figure 3.2 – Execution of trigonometric and logarithmic operations.
    \item Figure 3.3 – Error handling and result display for invalid inputs.
\end{itemize}

Each interface element was designed to contribute to user accessibility, operational clarity, and accurate computation.

\bigskip
\noindent
\textbf{3.3 Discussions and Observations}

Throughout the development process, the following key observations were made:
\begin{itemize}
    \item The separation of structure (HTML), style (CSS), and logic (JavaScript) proved essential for organized and maintainable code.
    \item Implementing trigonometric functions in degrees required custom conversions from degrees to radians, reinforcing understanding of mathematical transformations in JavaScript.
    \item Event-driven programming effectively handled user interaction and real-time display updates.
    \item While the calculator operates entirely on the client side, future enhancements could include persistent history logging or advanced symbolic computation.
\end{itemize}

Overall, the project demonstrated how foundational web technologies can be combined to create a fully functional and aesthetically refined application.

\bigskip
\noindent
\textbf{3.4 Summary}

This chapter presented the implementation results and observations for the Scientific Calculator web application. The calculator successfully performs both basic and scientific computations using a clean, interactive interface. The discussion highlighted the key aspects of design, functionality, and usability achieved during the development phase. The results affirm the effectiveness of using HTML, CSS, and JavaScript as introductory tools for practical web application development.


\newpage
\begin{center}
    \Large\textbf{CHAPTER 5}\\[0.3cm]
    \large\textbf{CONCLUSION}
\end{center}

The four-week industrial training program in Web Development provided a structured and foundational introduction to modern web technologies. The sessions offered a systematic progression from fundamental concepts of HTML and CSS to the basic understanding of JavaScript, equipping trainees with the essential knowledge required to design and implement static and interactive web pages. 

Throughout the training, emphasis was placed on practical implementation and conceptual clarity. The exposure to real development environments, coding standards, and responsive design principles contributed to developing a strong groundwork for future learning. While the scope of the program was introductory, the clarity achieved in basic front-end development concepts created a robust base for deeper exploration into dynamic and full-stack web application development.

The final project—a scientific web calculator—served as a synthesis of the acquired concepts. It demonstrated the integration of HTML for structure, CSS for layout and visual presentation, and JavaScript for interactivity. Although modest in complexity, the project represented an important step toward understanding client-side logic and the functional potential of web applications.

In conclusion, the industrial training successfully fulfilled its objective of introducing the core principles of web development to beginners. It provided not only theoretical comprehension but also practical competence, instilling confidence to independently pursue more advanced technologies and frameworks. The experience has laid a durable foundation for continued learning and professional growth in the evolving domain of web technologies.


% ------------------------------------------------
% REFERENCES
% ------------------------------------------------
\newpage
\begin{center}
    \Large\textbf{REFERENCES}
\end{center}

\begin{enumerate}
    \item W3C (World Wide Web Consortium). \textit{HTML Living Standard}. Available at: \url{https://html.spec.whatwg.org/} [Accessed July 2025].

    \item Mozilla Developer Network (MDN). \textit{HTML Documentation}. Available at: \url{https://developer.mozilla.org/en-US/docs/Web/HTML} [Accessed July 2025].

    \item Mozilla Developer Network (MDN). \textit{CSS: Cascading Style Sheets Documentation}. Available at: \url{https://developer.mozilla.org/en-US/docs/Web/CSS} [Accessed July 2025].

    \item Mozilla Developer Network (MDN). \textit{JavaScript Reference and Guide}. Available at: \url{https://developer.mozilla.org/en-US/docs/Web/JavaScript} [Accessed July 2025].

    \item Duckett, Jon. \textit{HTML and CSS: Design and Build Websites}. John Wiley and Sons, 2011.

    \item Duckett, Jon. \textit{JavaScript and JQuery: Interactive Front-End Web Development}. John Wiley and Sons, 2014.

    \item Robbins, Jennifer Niederst. \textit{Learning Web Design: A Beginner’s Guide to HTML, CSS, JavaScript, and Web Graphics}. 5th Edition, O’Reilly Media, 2018.

    \item Keith, Jeremy. \textit{HTML5 for Web Designers}. A Book Apart, 2010.

    \item Meyer, Eric A. \textit{CSS: The Definitive Guide}. 4th Edition, O’Reilly Media, 2017.

    \item Crockford, Douglas. \textit{JavaScript: The Good Parts}. O’Reilly Media, 2008.

    \item Flanagan, David. \textit{JavaScript: The Definitive Guide}. 7th Edition, O’Reilly Media, 2020.

    \item Bootstrap Team. \textit{Bootstrap Documentation}. Available at: \url{https://getbootstrap.com/docs/5.3/getting-started/introduction/} [Accessed July 2025].

    \item Visual Studio Code Documentation. \textit{User Guide and Extensions}. Microsoft Corporation. Available at: \url{https://code.visualstudio.com/docs} [Accessed July 2025].

    \item GitHub Documentation. \textit{Getting Started with Git and GitHub}. Available at: \url{https://docs.github.com/en/get-started} [Accessed July 2025].

    \item Khan Academy. \textit{Intro to HTML/CSS: Making Webpages}. Available at: \url{https://www.khanacademy.org/computing/computer-programming/html-css} [Accessed July 2025].

    \item freeCodeCamp. \textit{Responsive Web Design Certification}. Available at: \url{https://www.freecodecamp.org/learn/} [Accessed July 2025].

    \item W3Schools. \textit{HTML, CSS, and JavaScript Tutorials}. Available at: \url{https://www.w3schools.com/} [Accessed July 2025].

    \item Academic Advancement of Information Technology (A2IT), Mohali. \textit{Training Modules and Lecture Material on Web Development}, 2025.

    \item Guru Nanak Dev Engineering College, Ludhiana. \textit{Industrial Training Guidelines and Evaluation Criteria}, Department of Computer Science and Engineering, 2025.
\end{enumerate}

% ------------------------------------------------
% APPENDIX
% ------------------------------------------------
\appendix
\chapter*{APPENDIX}
(Include any additional material, code, or data here.)

\end{document}
